  \subsubsection{小数部分を残し,区間内の回転として読む}\label{節：ガウス型の小数部分と回転}
    \,$\{x\}=x-\lfloor x\rfloor$は$x$の小数部分を表し,\,$0\le\{x\}<1$である.
    例えば,\,$\{-0.2\}=0.8$となる.
    ある量を加えてから整数部分を捨てる操作は,値を$[0,1)$へ戻す操作である.
    区間の両端をつないで円周と見れば,毎回同じ割合だけ進む回転として解釈できる.

    \textbf{小数部分を追うと,整数の境界を越えても運動の規則が残る}.
    途中で何回境界を越えたかは,\bookterm{floor-function}{床関数}で記録できる.
    有理数の割合なら繰り返しが生じるが,無理数の場合については,周期があるかを別に証明しよう.

    \noindent \bookblocklink{example-gauss-role-9}{problem}{answer}{【例題】}\par\nobreak
    \begin{enumerate}[label=(\arabic*),parsep=3pt,beginpenalty=10000,format=\bookitemlink{example-gauss-role-9}{problem}{answer}]
      \item \,$\alpha=\sqrt2-1$とし,\bookterm{sequence}{数列}を次のように定める.
            \[
              a_0=0,\qquad
              a_{n+1}=a_n+\alpha-\lfloor a_n+\alpha\rfloor
              \quad(n\ge0).
            \]
            \bookterm{general}{一般項}を求め,異なる添字の項が同じ値になるか調べよ.
    \end{enumerate}

    \noindent \bookblocklink{example-gauss-role-9}{hint}{problem}{【方針】}\par\nobreak
    \begin{enumerate}[label=(\arabic*),parsep=3pt,beginpenalty=10000,format=\bookitemlink{example-gauss-role-9}{hint}{problem}]
      \item 整数部分を捨てると小数部分だけが残ることから,\,$n\alpha$の小数部分を予想します.
            一致する二項があれば,添字の差と$\alpha$の積が整数になることを使いましょう.
    \end{enumerate}

    \noindent \bookblocklink{example-gauss-role-9}{answer}{problem}{【解答】}\par\nobreak
    \begin{enumerate}[label=(\arabic*),parsep=3pt,beginpenalty=10000,format=\bookitemlink{example-gauss-role-9}{answer}{problem}]
      \item \[
              a_n=n\alpha-\lfloor n\alpha\rfloor\quad(n\ge0).
            \]
            異なる添字の項は同じ値にならず,周期的ではない.
    \end{enumerate}

    \noindent \bookblocklink{example-gauss-role-9}{solution}{problem}{【解法】}\par\nobreak
    \begin{enumerate}[label=(\arabic*),parsep=3pt,beginpenalty=10000,format=\bookitemlink{example-gauss-role-9}{solution}{problem}]
      \item \,$n=0$では予想した式は$a_0=0$を与える.
            \,$a_n=n\alpha-\lfloor n\alpha\rfloor$とすると,整数との和の性質から,
            \begin{align*}
              \lfloor a_n+\alpha\rfloor
                &=\lfloor(n+1)\alpha\rfloor-\lfloor n\alpha\rfloor,\\
              a_{n+1}
                &=(n+1)\alpha-\lfloor(n+1)\alpha\rfloor.
            \end{align*}
            よって帰納法により\bookterm{general}{一般項}が示された.すべての項は$[0,1)$に属する.

            \,$a_{n+p}=a_n\ (p\ge1)$とすると,
            \[
              p\alpha=\lfloor(n+p)\alpha\rfloor-\lfloor n\alpha\rfloor
            \]
            が整数になり,\,$\alpha$が有理数となる.これは$\alpha=\sqrt2-1$が無理数であることに矛盾する.
            従って異なる添字の項は一致しない.

            同じ\bookterm{recurrence}{漸化式}で$\alpha$が有理数$\frac pq\ (0<p<q)$であり,\,$p,q$が互いに素なら,
            最小の正の周期は$q$となる.\,$q\alpha$が整数なので$q$項後に戻り,
            戻るまでの項数$T$は$T\alpha$が整数となる条件から$q$の倍数だからである.
            回転の見方と,有理数・無理数という整数問題の考え方が結びついている.
    \end{enumerate}
